\VerdierPhysicalPageAnchor{0079}%
\VerdierSemanticAnchor{ch1-2-5-7-morphisms}
\noindent\hspace*{2em}--\ Let $X$ and $Y$ be two objects of $K(C)$, and let
$f,g\colon X\rightrightarrows Y$ be two morphisms of
$\operatorname{comp}(C)$. The relation ``$f$ and $g$ are homotopic'' is an
equivalence relation on $\operatorname{Hom}_{\operatorname{comp}(C)}(X,Y)$,
and $\operatorname{Hom}_{K(C)}(X,Y)$ is defined to be the quotient of
$\operatorname{Hom}_{\operatorname{comp}(C)}(X,Y)$ by this equivalence
relation.

\VerdierSemanticAnchor{ch1-2-5-7-composition}
\noindent\hspace*{2em}--\ It follows immediately from (2.5.5) that the law of
composition on the morphisms of $\operatorname{comp}(C)$ passes to the
quotient and defines a law of composition on the morphisms of $K(C)$.

\VerdierSemanticAnchor{ch1-2-5-7-additive-category}
It follows immediately from (2.5.5) that we have indeed thereby defined a
category and that this category is additive.

\VerdierSemanticAnchor{ch1-2-5-7-quotient-functor}
The canonical functor $\operatorname{comp}(C)\xrightarrow{Q}K(C)$, which is
the identity on objects and associates to every morphism $f$ of
$\operatorname{comp}(C)$ its equivalence class $\dot f$, is additive.

\VerdierSemanticAnchor{ch1-2-5-7-homotopy-class}
The morphism $\dot f$ of $K(C)$, the image of the morphism $f$ of
$\operatorname{comp}(C)$ under the canonical functor, is called the
\emph{class of $f$ modulo homotopy}.

\VerdierSemanticAnchor{ch1-2-5-7-name}
The category $K(C)$ is called the \emph{category of complexes of $C$ up to
homotopy}.

\VerdierSemanticAnchor{ch1-2-5-8}
\noindent\textbf{2.5.8.}\quad
The canonical functor $\operatorname{comp}(C)\xrightarrow{Q}K(C)$ has the
following property: every functor from $K(C)$ to a category $A$ yields, by
composition with $Q$, a functor from $\operatorname{comp}(C)$ to $A$; whence,
for every category $A$, a functor:

\VerdierSemanticAnchor{ch1-2-5-8-hom-functor}
\begin{equation*}
  \mathcal H\!om(Q,A)\colon
  \mathcal H\!om(K(C),A)
  \longrightarrow
  \mathcal H\!om(\operatorname{comp}(C),A)
\end{equation*}
($\mathcal H\!om$ denotes the category of functors).

\VerdierSemanticAnchor{ch1-2-5-8-fully-faithful}
This functor is fully faithful, injective on objects, and its image is the full
subcategory of:

\VerdierSemanticAnchor{ch1-2-5-8-essential-image}
\begin{equation*}
  \mathcal H\!om(\operatorname{comp}(C),A)
\end{equation*}
defined by the functors from $\operatorname{comp}(C)$ to $A$ that send every
pair of homotopic morphisms of $\operatorname{comp}(C)$ to equal morphisms.

\VerdierSemanticAnchor{ch1-2-5-8-uniqueness}
The category $K(C)$, equipped with the canonical functor
$\operatorname{comp}(C)\xrightarrow{Q}K(C)$, is determined up to unique
isomorphism by this property.

\VerdierSemanticAnchor{ch1-2-5-9}
\noindent\textbf{2.5.9.}\quad
It follows immediately from (2.5.5) and (2.5.8) that there exists one and only
one translation functor on $K(C)$ such that the canonical functor

\VerdierSemanticAnchor{ch1-2-5-9-translation}
\begin{equation*}
  \operatorname{comp}(C)\longrightarrow K(C)
\end{equation*}
commutes (strictly) with translations.

\VerdierSemanticAnchor{ch1-2-5-9-strict-z-category}
This translation functor is an automorphism of $K(C)$ and consequently
defines on $K(C)$ a strict $\mathbb Z$-category structure.

\VerdierSemanticAnchor{ch1-2-5-9-strict-z-functor}
The canonical functor $\operatorname{comp}(C)\longrightarrow K(C)$ is a
strict $\mathbb Z$-functor.
